% Pista pro-24s-q4 · Probabilitat
% Procedència: PAU Catalunya 2024 · Convocatòria de setembre · Sèrie 3 · Qüestió 4
\documentclass[11pt,a4paper]{article}
\usepackage{pista}
\pistainfo{Catalunya 2024}{Convocatòria de setembre}{Sèrie 3}

\begin{document}

\section*{\textcolor{blauPista}{Qüestió 4}}

\noindent S'estima que el $20\,\%$ dels habitants d'una regió pateix algun tipus d'arrítmia. Per a diagnosticar-la, hi ha la possibilitat de col·locar al pacient un monitor Holter, que detecta l'arrítmia en un $95\,\%$ dels casos de persones que la pateixen, però que també dona falsos positius, per motius elèctrics, en persones que no pateixen arrítmies en un $0{,}5\,\%$ dels casos.

\subsection*{\textit{a)} \normalfont Si escollim 4 persones a l'atzar, quina és la probabilitat que almenys una d'elles pateixi arrítmies? \hfill\textnormal{[0,75~punts]}}

\pista{Aquí tens una variable binomial $X \sim B(4, 0{,}2)$, on $0{,}2$ és la probabilitat que una persona qualsevol pateixi arrítmies. Molt sovint, quan et demanen treballar amb el succés ``almenys una'', és més còmode treballar amb el succés contrari, que en aquest cas, seria  ``cap d'elles''. Per tant, has de calcular $q= Prob(X = 0)$ i, per obtenir la resposta final $p$, recorda que hauràs de "tornar al camí original", és a dir, que la resposta serà $p = 1 - q$.}

\subsection*{\textit{b)} \normalfont Quina és la probabilitat que una persona escollida a l'atzar obtingui un diagnòstic positiu d'arrítmia? \hfill\textnormal{[0,75~punts]}}

\pista{Un diagrama d'arbre t'ajudarà a visualitzar els dos camins que porten a un diagnòstic positiu: o bé la persona pateix arrítmia i el Holter la detecta, o bé no la pateix però el Holter dona un fals positiu. Aplica el teorema de la probabilitat total.}

\subsection*{\textit{c)} \normalfont Si una persona obté un diagnòstic negatiu a la prova del Holter, quina és la probabilitat que realment pateixi arrítmies? \hfill\textnormal{[1~punt]}}

\pista{Aquest és un exemple clar on has d'aplicar el Teorema de Bayes, perquè et demanen una probabilitat condicionada en sentit \emph{invers} respecte de la informació que et donen: tu coneixes la probabilitat de ``donar negatiu \emph{sabent que} pateix arrítmia'' (que és el complementari del $95\,\%$ de detecció), però et demanen la probabilitat de ``patir arrítmia \emph{sabent que} ha donat negatiu''. Hi ha un detall important que sovint passa amb aquests problemes, i és que els apartats estan relacionats. M'explico: la probabilitat de ``donar negatiu'' la pots obtenir gairebé immediatament, pensant que és el succés contrari al de l'apartat b). Per tant, serà $1 - $ "prob. positiu".}

\end{document}
